% ---------------------------------------------------------------------------
% 6 Optimization -- the systems core of the paper
% Source: evidence/optimizations.yaml (authoritative), docs/hybrid/t123_walkthrough.md §7, §9
% Discipline: the operation schedule is identical at every step. Say so up front; it is what
% makes the speedup attributable to systems work rather than to a cheaper circuit.
%
% Tone: this section is the good news and should read that way. Abandoned attempts appear only
% where they explain a retained design choice (the clone, and the sequential client boundary),
% one or two sentences each, in the subsection they explain. No catalogue of failures.
% Only timing artifacts that satisfy the manuscript evidence rule may appear in the paper.
% ---------------------------------------------------------------------------

\section{Host-side systems work at fixed circuit}
\label{sec:optimization}

Encoding each distinct weight diagonal once removes redundant host work but creates a memory
lifetime problem. An implementation that retained every encoded template was terminated by the
operating system at $61{,}863$~MiB resident, when the MLP began building templates on top of those
cached by earlier stages. Bounding the cache to one live stage makes reuse executable within the
tested host-memory envelope. Peak host memory is then determined by the largest stage, not by every
encoded weight in the block. The remaining changes govern when host data are prepared, how long
they remain resident, and where CPU work executes.

Every change in this section holds the encrypted circuit fixed: the retained implementation executes
the same checked operation schedule before and after each change, with packing, multiplicative
depth, boundary order, and the frozen plaintext reference unchanged, and the run-time checks of
Section~\ref{sec:protocol} abort on any deviation. No retained change removes a dense transform,
reduces the number of attention tiles, moves linear algebra to the client, or relaxes the acceptance
criterion. The section therefore reports verified changes in executed work and memory but no
elapsed-time comparison, because no dedicated-node measurement of the pre-optimization configuration
exists.

\subsection{Host-side encoding redundancy}
The baseline rebuilt and encoded every packed diagonal for every dense product. Its $156$ products
each required $1{,}024$ diagonal plaintexts, producing $159{,}744$ encoding operations per block,
and about $90\%$ of that work was byte-identical: the block has $12$ weight matrices, each used
across all $13$ token groups. The GPU could not begin a multiplication until the host had prepared
its plaintexts.

In the baby-step giant-step transform, a plaintext diagonal is not merely a view of the original
matrix: it must be packed, scaled, and encoded for the ciphertext level at which multiplication will
occur. Repeating that conversion for every token group made a model-invariant representation step
dominate the call path. The optimization target is therefore the lifetime and placement of plaintext
preparation, while the encrypted products and accumulations remain fixed.

\subsection{Encode once, clone per use}
The retained cache encodes each distinct weight diagonal into a pristine host template. A
multiplication receives a lightweight clone that shares the encoded host data by reference while
carrying fresh device state; the clone is loaded and evicted for that operation. This removes
$90\%$ of the $159{,}744$ encoding operations without keeping a plaintext object resident on the
GPU between multiplications.

Cloning is required by the evaluated backend's object semantics. Its plaintext-load path returns
early when an object is already marked as loaded, so directly reusing a template across
multiplicative levels can pair a new ciphertext level with stale residue-number-system limbs.
Fresh device state on each clone forces a load at the level required by the current operation. An
isolated GPU micro-test checked clone-versus-fresh-encode equality, cross-level reuse, stable device
memory under repeated multiplication, and equality between parallel and serial encoding before the
cache was applied to the complete block.

Templates are indexed by both weight matrix and multiplicative level, which preserves the backend's
level contract while still sharing work among token groups that reach the same matrix at the same
level. The cache therefore stores reusable encoded data, not a ciphertext or a mutable GPU object,
and leaves homomorphic multiplication itself unchanged.

\subsection{Bounding host memory}
Removing repeated encoding exchanges computation for host memory. An early cache retained
templates from every completed stage, and the process did not survive the MLP. Two exact changes
bound the cache. First, a template is encoded directly at its multiplicative level of use, so it
carries only the limbs needed by the receiving ciphertext; this changes no arithmetic, because the
backend would otherwise drop a level-$0$ plaintext to the same level before multiplication. Second,
the implementation flushes the cache after query/key/value projection and after the attention
projection, which are the final use sites for their respective weight sets: no weight matrix recurs
across either boundary, so retaining those templates cannot produce another hit.

The lifetime policy follows the block topology. Within a stage, the group-major loop revisits its
weights as it advances through token groups, so early eviction would discard real reuse; across
stages, the next operation consumes a different set of weights, so eviction there releases memory
without creating a later re-encoding.

The target-process resident-set trace follows this policy, as
Figure~\ref{fig:systems}(d) shows. Query/key/value projection reaches
$25{,}530$~MiB. Attention begins after a flush and uses $12{,}434$~MiB. The MLP holds eight
matrices and reaches the block peak of $48{,}610$~MiB. The largest live stage, rather than every
encoded weight in the block, determines peak host memory.

\subsection{Placement and parallelism}
Plaintext preparation also depends on where host work runs. The retained implementation binds
threads to physical cores and uses NUMA-local allocation so encoding threads stay close to the
memory they allocate and to buffers staged for the GPU; decryption, exact nonlinear evaluation, and
re-encryption use the same host resources. Template construction is then batched across the
available $32$ cores with dynamic scheduling, since each diagonal can be encoded independently
without changing the order or result of encrypted multiplications; parallel and serial construction
were checked for equal outputs before the parallel path was retained. Affinity controls placement
and batching exposes independent diagonals to the cores, so together they make plaintext preparation
a core-local parallel batch under the same cryptographic parameters. Their elapsed-time effect
awaits a sourced paired measurement.

\begin{table*}[!tb]\centering\small
\caption{Host-side changes at fixed circuit. Every row leaves the operation schedule, packing,
multiplicative depth, and frozen plaintext reference unchanged. Effects are verified changes in
executed work or in memory; the table carries no elapsed-time column, because no dedicated-node
measurement of the pre-optimization configuration exists.}
\label{tab:optimization}
\begin{tabularx}{\linewidth}{Y l l}
\toprule
Change & Verified effect & Status \\
\midrule
Re-encode every plaintext per multiplication & $159{,}744$ encodings & baseline \\
Thread affinity, NUMA-local allocation & host placement changed & retained \\
Encode once, clone per use & removes $90\%$ of encodings & retained \\
Parallel batched encoding & equal to serial encoding & retained \\
Encode at use level, flush between stages & peak $48{,}610$~MiB RSS & enabling \\
\midrule
\textbf{Net} & \textbf{checked operation schedule unchanged} & \\
\bottomrule
\end{tabularx}
\end{table*}

% No waterfall figure: it needs two supported timing endpoints, and the pre-optimization
% baseline has no committed dedicated-node measurement. The generator is retained in
% scripts/figures.py but is not registered in FIGURES.

\subsection{Remaining optimization targets}
Mask plaintexts are still encoded afresh at approximately $17{,}400$ uses per block. Head masks,
score-isolation masks, lane masks, and copy-activity masks all have fixed values under the declared
layout, so host templates could remove the same class of repeated work already eliminated for
weight diagonals. This target remains unimplemented.

A second opportunity lies in the four packed activation copies. They currently carry duplicates
through a dense transform. An analyzed layout would apply distinct transforms to those copies,
placing query, key, and value projections together and likewise combining the MLP expansion
chunks. At fixed model arithmetic, the projected schedule reduces dense products from $156$ to
approximately $52$ and dense ciphertext--plaintext multiplications from $159{,}744$ to
approximately $53{,}248$; the total ciphertext--plaintext count falls by roughly $60\%$ before mask
and copy-repair overhead. This is a derived projection for an unimplemented layout, and it requires
an exact reconstruction proof before implementation.

Client work has a separate opportunity. Attention dominates the boundary schedule
(Section~\ref{sec:protocol}), and each attention tile's cryptographic work is independent.
Concurrent client boundaries were unsafe under the evaluated shared crypto context, so the retained
implementation executes them sequentially. A thread-safe context would expose this independent work
to concurrency.
